\tikzset{
  hw/.style={draw, rounded corners=1pt, align=center, font=\small, inner sep=4pt,
    minimum height=9mm},
  hwmem/.style={hw, draw=violet!65!black, fill=violet!8},
  hwreg/.style={hw, draw=blue!60!black, fill=blue!9},
  hwop/.style={hw, draw=orange!65!black, fill=orange!13},
  hwctl/.style={hw, draw=black!65, fill=black!5},
  hwmux/.style={hw, trapezium, trapezium left angle=70, trapezium right angle=70,
    shape border rotate=270, draw=teal!65!black, fill=teal!12},
  hwdata/.style={-{Stealth[length=2mm]}, thick, draw=black!78},
  hwctrl/.style={-{Stealth[length=2mm]}, dashed, draw=black!62},
  hwlabel/.style={font=\scriptsize, align=center, inner sep=2pt}
}

\clearpage
\subsection{Jalur masukan dan pemilihan sumber piksel}
Diagram rangkaian berikut memakai Baseline A dengan satu jalur MAC Conv1
dan satu jalur MAC FC. Kotak biru menunjukkan register yang menyimpan nilai
pada tepi naik clock. Ungu menunjukkan memori, hijau menunjukkan MUX, dan
oranye menunjukkan logika kombinasi. Garis putus-putus membawa sinyal kontrol.
Simbol register pada keluaran memori menunjukkan latensi baca satu clock.
Vivado dapat memetakan fungsi register itu ke dalam primitif memori.

\begin{figure}[H]
\centering
\begin{tikzpicture}
\node[hwctl,text width=24mm] (uart) at (1.4,4.5) {RsRx\\UART RX};
\node[hwctl,text width=25mm] (frame) at (4.65,4.5) {\texttt{frame\_ctrl}\\data, addr, WE};
\node[hwmem,text width=27mm] (live) at (8.35,4.5) {LIVE RAM\\\(256\times8\) bit};
\node[hwreg,text width=19mm] (lp) at (11.1,4.5) {\texttt{live\_pixel}\\8 bit};
\node[hwmem,text width=24mm] (stored) at (8.35,2.35) {Stored ROM\\\(4096\times8\) bit};
\node[hwreg,text width=23mm] (sp) at (11.35,2.35) {\texttt{stored\_pixel}\\8 bit};
\node[hwmux] (mux) at (13.55,3.42) {MUX};
\node[hwctl,text width=14mm] (conv) at (15.3,3.42) {Conv1};
\node[hwctl,text width=24mm] (sw) at (1.4,2.35) {SW3--SW0\\sinkronisasi};
\node[hwreg,text width=33mm] (slot) at (4.65,2.35) {\texttt{stored\_selection}\\4 bit};
\node[hwreg,text width=25mm] (mode) at (13.55,0.48) {\texttt{active\_uart}\\1 bit};
\node[hwctl,text width=28mm] (arb) at (4.65,0.48) {Arbiter start\\\texttt{top\_basys3}};
\draw[hwdata] (uart) -- (frame);
\draw[hwdata] (frame) -- (live);
\draw[hwdata] (live) -- (lp);
\draw[hwdata] (stored) -- (sp);
\draw[hwdata] (lp.east) -- ++(3mm,0) |- ([yshift=2.2mm]mux.west);
\draw[hwdata] (sp.east) -- ++(1.8mm,0) |- ([yshift=-2.2mm]mux.west);
\draw[hwdata] (mux) -- (conv);
\draw[hwdata] (sw) -- (slot);
\draw[hwdata] (slot) -- node[hwlabel,above]{slot} (stored);
\draw[hwctrl] (arb.north) -- node[hwlabel,right]{start stored} (slot.south);
\draw[hwctrl] (arb) -- (mode);
\draw[hwctrl] (mode) -- (mux);
\draw[hwctrl] (frame.south) -- node[hwlabel,right]{checksum\\lolos} (4.65,3.3)
  -| (2.93,0.85) -- ([yshift=3mm]arb.west);
\node[hwlabel,anchor=west] at (0.12,0.6) {BTNC saat\\SW15 = 0};
\draw[hwctrl] (2.0,0.48) -- (arb);
\end{tikzpicture}
\caption{Jalur masukan RTL. MUX memilih piksel LIVE saat \texttt{active\_uart}=1 dan piksel stored saat nilainya 0. Kedua memori membaca alamat dari Conv1.}
\label{fig:hw-input}
\end{figure}

Host mengirim piksel \(16\times16\) yang telah dipraproses. UART RX
merakit byte serial. \texttt{frame\_ctrl} menulis byte piksel ke LIVE
RAM selama penerimaan dan meminta inferensi setelah checksum lolos.
Untuk citra stored, BTNC memicu start ketika SW15 rendah. Arbiter
menyimpan pilihan slot dan sumber saat menerima start sehingga perubahan
switch di tengah inferensi tidak mengubah piksel yang dibaca.

Conv1 mengirim \texttt{pix\_addr} 8 bit ke kedua memori. LIVE RAM
membaca alamat tersebut langsung. Stored ROM membaca gabungan
\(\{\texttt{stored\_selection},\texttt{pix\_addr}\}\) yang panjangnya
12 bit. Jadi tiap slot stored memiliki 256 alamat berurutan. Pada clock
berikutnya, MUX meneruskan salah satu keluaran baca 8 bit ke Conv1.
Setelah Conv1 selesai, \texttt{cnn\_controller} memulai classifier.
Classifier mengurutkan MaxPool, FC, dan argmax. Top mengarahkan ID
kelas ke LED dan seven-segment. Untuk transaksi LIVE, top juga mengirim
ID, skor, dan counter siklus melalui UART.

\clearpage
\subsection{Konvolusi dan requantization}
\texttt{conv\_engine} menyimpan indeks filter \(f\), posisi keluaran
\((r,c)\), dan posisi tap \((k_r,k_c)\) pada register pencacah. Pada setiap
langkah serial, logika alamat membentuk
\(a_{\rm pix}=16(r+k_r)+(c+k_c)\) dan
\(a_w=9f+3k_r+k_c\). Untuk filter yang sama, sembilan bobot digunakan
ulang ketika jendela bergeser. Setelah \(14\times14\) posisi selesai,
pencacah maju ke filter berikutnya.

\begin{figure}[H]
\centering
\begin{tikzpicture}
\node[hwreg,text width=24mm] (cnt) at (1.5,2.65) {Pencacah\\\(f,r,c,k_r,k_c\)};
\node[hwmem,text width=25mm] (pix) at (4.65,4.4) {Port piksel\\baca 1 clock};
\node[hwmem,text width=25mm] (rom) at (4.65,2.65) {ROM bobot\\\(72\times8\) bit};
\node[hwmem,text width=25mm] (bias) at (4.65,1.2) {ROM bias\\\(8\times24\) bit};
\node[hwop] (mul) at (7.1,3.45) {\(\times\)};
\node[hwreg,text width=16mm] (p) at (9.0,3.45) {Produk\\signed17};
\node[hwmux] (mux) at (9.0,1.2) {MUX};
\node[hwop,text width=16mm] (sum) at (11.35,2.55) {\(+\)\\signed25};
\node[hwreg,text width=18mm] (acc) at (13.9,2.55) {\texttt{acc}\\signed24};
\node[hwctl,text width=22mm] (out) at (13.9,0.45) {Ke requant\\pada tap ke-9};
\draw[hwdata] (cnt.east) -- ++(3mm,0) |- (pix.west);
\draw[hwdata] (cnt) -- (rom);
\draw[hwdata] (cnt.east) -- ++(3mm,0) |- (bias.west);
\draw[hwdata] (pix.east) -| (mul.north);
\draw[hwdata] (rom.east) -| (mul.south);
\draw[hwdata] (mul) -- (p);
\draw[hwdata] (p.east) -| (sum.north);
\draw[hwdata] (bias.east) -- node[hwlabel,above]{bias} (mux.west);
\draw[hwdata] (mux.east) -| (sum.south);
\draw[hwdata] (sum) -- (acc);
\draw[hwdata] (13.9,1.68) -| (12.2,0.05)
  -- node[hwlabel,below]{umpan balik \texttt{acc}} (9.0,0.05) -- (mux.south);
\draw[hwdata] (acc) -- (out);
\fill (13.9,1.68) circle (1.4pt);
\end{tikzpicture}
\caption{Datapath Conv1 serial. Satu pengali dipakai untuk sembilan tap per keluaran. MUX memilih bias filter pada tap pertama, lalu jumlah parsial dari register accumulator.}
\label{fig:hw-conv}
\end{figure}

Piksel \texttt{uint8} diperluas dengan nol sebelum dikalikan dengan bobot
\texttt{int8}. Register produk menyimpan hasil signed17. Bias dibaca
menggunakan nomor filter yang sudah diregister agar tiba bersama produk.
Penjumlah 25 bit memeriksa overflow dari perbedaan dua bit teratas.
Register \texttt{acc} menyimpan 24 bit terbawah. Sinyal valid, penanda
tap pertama/terakhir, dan alamat keluaran ikut berjalan sepanjang pipeline.
Satu keluaran lengkap tersedia setelah sembilan produk. Setelah pipeline
terisi, datapath menerima satu tap per clock.

\begin{figure}[H]
\centering
\begin{tikzpicture}
\node[hwreg] (acc) at (0.68,0) {Acc\\signed24};
\node[hwmux,text width=16mm] (relu) at (2.7,0) {ReLU\\\(0\) atau acc};
\node[hwop,text width=15mm] (mul) at (4.66,0) {\(\times M\)\\43980};
\node[hwreg,text width=15mm] (p) at (6.65,0) {Produk\\39 bit};
\node[hwop,text width=18mm] (shift) at (9.05,0) {\(+2^{24}\)\\geser 25 bit};
\node[hwmux,text width=15mm] (clip) at (11.4,0) {Clamp\\ke 127};
\node[hwreg] (q) at (13.3,0) {Aktivasi\\8 bit};
\node[hwmem] (ram) at (15.04,0) {RAM\\fitur};
\draw[hwdata] (acc)--(relu); \draw[hwdata] (relu)--(mul);
\draw[hwdata] (mul)--(p); \draw[hwdata] (p)--(shift);
\draw[hwdata] (shift)--(clip); \draw[hwdata] (clip)--(q);
\draw[hwdata] (q)--(ram);
\end{tikzpicture}
\caption{Requantization dua tahap pada A. Bit tanda memilih nol untuk ReLU. Setelah perkalian terdaftar, logika pembulatan, pergeseran tetap, dan saturasi membentuk aktivasi 0--127.}
\label{fig:hw-requant}
\end{figure}

Penjumlahan konstanta dan pergeseran pada gambar mewakili operasi
\((P M+2^{24})\mathbin{\gg}25\) dengan \(P=\max(0,\texttt{acc})\).
Karena jumlah geser konstan, logika dapat mengambil irisan bit.
Pembanding terhadap 127 dan MUX menjalankan saturasi. Keluaran
terdaftar membawa alamat \(196f+14r+c\) dan write-enable ke port tulis
RAM fitur \(1568\times8\) bit.

\clearpage
\subsection{MaxPool dan fully connected}
MaxPool membaca RAM fitur satu nilai per tap dengan FSM
\texttt{ISSUE}/\texttt{SAMPLE}. Untuk channel \(f\) dan posisi pooling
\((r_p,c_p)\), alamat tap \((t_r,t_c)\) adalah
\(196f+14(2r_p+t_r)+2c_p+t_c\). Empat alamat itu membentuk blok
\(2\times2\) pada satu channel. Stride 2 berasal dari logika alamat.

\begin{figure}[H]
\centering
\begin{tikzpicture}
\node[hwmem,text width=20mm] (mem) at (1.15,1.3) {RAM fitur\\baca 1 clock};
\node[hwop,text width=28mm] (cmp) at (4.4,2.45) {Pilih \(x\) jika tap 0\\atau \(x>\texttt{maximum}\)};
\node[hwmux] (mux) at (7.0,1.3) {MUX};
\node[hwreg,text width=23mm] (max) at (9.8,1.3) {\texttt{maximum}\\uint8};
\node[hwreg,text width=22mm] (out) at (9.8,-0.1) {\texttt{pool\_data}\\uint8};
\node[hwmem,text width=26mm] (pool) at (13.55,-0.1) {RAM pooling\\\(392\times8\) bit};
\draw[hwdata] (mem.east) -- node[pos=.25,hwlabel,below]{\(x\)} (mux.west);
\draw[hwdata] (3.0,1.3) |- (cmp.west);
\draw[hwctrl] (cmp.east) -| (mux.north);
\draw[hwdata] (mux) -- (max);
\draw[hwdata] (max.north) -- (9.8,3.38) -| (cmp.north);
\draw[hwdata] (max.east) -- (15.28,1.3) -- (15.28,-0.96)
  -- node[hwlabel,below]{umpan balik} (7.0,-0.96) -- (mux.south);
\draw[hwdata] (8.25,1.3) |- (out.west);
\draw[hwdata] (out) -- (pool);
\fill (3.0,1.3) circle (1.4pt);
\fill (8.25,1.3) circle (1.4pt);
\end{tikzpicture}
\caption{Datapath MaxPool. Pada tap pertama, MUX mengabaikan hasil pembanding. Pada tap keempat, hasil maksimum juga masuk register tulis RAM pooling.}
\label{fig:hw-pool}
\end{figure}

Pembanding unsigned dan MUX digunakan kembali pada empat tap.
Saat tap terakhir, \texttt{pool\_we} dan alamat keluaran juga
diregister. RAM menangkap hasil pada tepi clock berikutnya. Keluaran
ditulis berurutan dalam CHW, sehingga flatten \(8\times7\times7\)
menjadi 392 fitur cukup dengan membaca alamat 0--391. Susunan data itu
langsung sesuai dengan urutan bobot FC.

\begin{figure}[H]
\centering
\begin{tikzpicture}
\node[hwmem,text width=25mm] (pool) at (1.5,3.12) {RAM pooling\\alamat \(i\)};
\node[hwmem,text width=25mm] (w) at (1.5,1.72) {ROM bobot FC\\alamat \(392k+i\)};
\node[hwmem,text width=25mm] (b) at (1.5,0.1) {ROM bias FC\\alamat \(k\)};
\node[hwop,text width=13mm] (mul) at (4.1,2.42) {\(\times\)\\signed17};
\node[hwmux] (mux) at (4.1,0.1) {MUX};
\node[hwop,text width=15mm] (sum) at (6.75,1.2) {\(+\)\\signed25};
\node[hwreg,text width=17mm] (acc) at (9.3,1.2) {\texttt{acc}\\signed24};
\node[hwreg,text width=26mm] (scores) at (11.5,2.95) {Register skor\\\(8\times24\) bit};
\node[hwctl,text width=18mm] (arg) at (14.65,2.95) {Argmax};
\draw[hwdata] (pool.east) -| (mul.north);
\draw[hwdata] (w.east) -| (mul.south);
\draw[hwdata] (b.east) -- (mux.west);
\draw[hwdata] (mul.east) -| (sum.north);
\draw[hwdata] (mux.east) -| (sum.south);
\draw[hwdata] (sum)--(acc);
\draw[hwdata] (acc.south) -- (9.3,-0.96) -| (4.1,-0.55) -- (mux.south);
\node[hwlabel,fill=white] at (6.75,-0.96) {umpan balik untuk \(i>0\)};
\draw[hwdata] (8.0,1.2) |- (8.0,2.95)
  -- node[hwlabel,above]{jika \(i=391\)} (scores.west);
\fill (8.0,1.2) circle (1.4pt);
\draw[hwdata] (scores) -- (arg);
\end{tikzpicture}
\caption{Datapath FC serial. Satu jalur mengolah fitur \(i=0,\ldots,391\) untuk kelas \(k\). Pada fitur terakhir, skor mengambil keluaran penjumlah pada clock yang sama dengan penulisan accumulator.}
\label{fig:hw-fc}
\end{figure}

Pada \texttt{ISSUE}, FC menyediakan alamat aktivasi dan bobot. Pada
\texttt{MAC}, data baca diproses oleh pengali dan penjumlah.
MUX memilih bias signed24 untuk \(i=0\) dan accumulator untuk
langkah berikutnya. Setelah 392 fitur, register skor kelas \(k\)
menerima \texttt{sum[23:0]} pada irisan \texttt{scores[24*k +: 24]}.
Pembanding dua bit teratas penjumlah mengawasi overflow di setiap MAC.
Pada varian FC paralel, tiap jalur mendapat aktivasi yang sama dan
membaca bank bobot serta bias untuk kelompok kelasnya sendiri.

\clearpage
\subsection{Argmax dan pemetaan ke FPGA}
Argmax memuat skor kelas 0 ke register \texttt{best}, lalu memeriksa
kelas 1--7 satu per satu. Pemilih irisan bus mengambil 24 bit skor
sesuai register indeks. Pembanding signed dan write-enable memperbarui
skor dan indeks terbaik hanya saat kandidat lebih besar. Jika dua skor
sama, indeks terkecil bertahan. Pada kandidat terakhir, logika keluaran
memilih indeks berdasarkan hasil pembanding saat itu.

\begin{figure}[H]
\centering
\begin{tikzpicture}
\node[hwreg,text width=24mm] (scores) at (1.4,2.35) {Register skor\\\(8\times24\) bit};
\node[hwmux,text width=18mm] (mux) at (4.2,2.35) {Pilih\\skor indeks};
\node[hwop,text width=22mm] (cmp) at (7.3,2.35) {Pembanding\\signed \(>\)};
\node[hwreg,text width=28mm] (best) at (11.1,2.35) {\texttt{best} +\\\texttt{best\_index}};
\node[hwreg,text width=22mm] (idx) at (4.2,0.25) {Pencacah\\indeks kelas};
\node[hwmux,text width=19mm] (last) at (11.1,0.25) {Kelas akhir\\pilih indeks};
\node[hwreg,text width=18mm] (class) at (14.0,0.25) {\texttt{class\_id}\\3 bit};
\draw[hwdata] (scores)--(mux);
\draw[hwdata] (idx)--(mux);
\draw[hwdata] (mux)--(cmp);
\draw[hwctrl] (cmp)-- node[hwlabel,above]{WE} (best);
\draw[hwdata] (best.north) -- (11.1,3.46) -| (cmp.north);
\draw[hwdata] (5.6,2.35) |- (10.15,1.48) -- ([xshift=-9.5mm]best.south);
\node[hwlabel,above] at (7.35,1.48) {skor kandidat};
\fill (5.6,2.35) circle (1.4pt);
\draw[hwdata] (idx) -- node[hwlabel,above,pos=.46]{indeks kandidat} (last);
\draw[hwdata] (best)--(last);
\draw[hwctrl] (8.9,2.35) |- (9.8,-1.18) -| (last.south);
\draw[hwdata] (last)--(class);
\end{tikzpicture}
\caption{Datapath argmax serial. Pemilih skor dan pembanding signed24 digunakan kembali. Keluaran 3 bit menyimpan indeks RTL kelas pemenang.}
\label{fig:hw-argmax}
\end{figure}

Diagram sebelumnya menjelaskan rangkaian pada tingkat RTL.
Pemetaan berikut diambil dari checkpoint pasca-routing lokal
\texttt{build/basys3/top\_basys3\_routed.dcp}, Vivado 2026.1 pada
XC7A35T. Nama lokasi berasal dari properti \texttt{LOC} sel netlist.
Vivado menggabungkan logika aritmetika ke DSP48E1 dan memakai
register internal pengali (\texttt{MREG=1}) pada Conv1 dan requant.
Netlist juga memuat 8 flip-flop \texttt{FDRE} untuk
\texttt{maximum} MaxPool dan 192 untuk register skor FC.

\begin{table}[H]
\centering\small
\caption{Pemetaan blok dari laporan hierarkis checkpoint lokal Baseline A.}
\begin{tabularx}{\linewidth}{llX}
\toprule
Blok & Primitif & Lokasi \\
\midrule
LIVE RAM & RAMB18E1 & \texttt{RAMB18\_X0Y10} \\
Stored ROM & RAMB36E1 & \texttt{RAMB36\_X0Y6} \\
Conv1 \texttt{acc\_sum} & DSP48E1 & \texttt{DSP48\_X0Y13} \\
Requant \texttt{rounded} & DSP48E1 & \texttt{DSP48\_X0Y14} \\
RAM fitur & RAMB18E1 & \texttt{RAMB18\_X0Y14} \\
RAM pooling & RAMB18E1 & \texttt{RAMB18\_X0Y16} \\
FC \texttt{sum} & DSP48E1 & \texttt{DSP48\_X0Y20} \\
\bottomrule
\end{tabularx}
\label{tab:hw-physical}
\end{table}

Checkpoint lokal memakai 1.104 LUT dan 932 FF. Instans
\texttt{u\_fc} memakai 541 LUT, 246 FF, satu DSP dan nol BRAM.
Nilai 1.106 LUT dan 934 FF pada tabel perbandingan model berasal dari
build PR~\#7. D42 menambah Conv2 dan penyimpanan aktivasi sebelum FC.
C-43 memakai ROM bobot FC di BRAM dan memperlebar keluaran menjadi
43 skor serta indeks kelas 6 bit \cite{repo}.
